\section{Error Handling}


\subsection{Precompile Error Semantics}

When a precompile returns an error, the EVM:
\begin{enumerate}[nosep]
  \item Consumes all gas provided to the call.
  \item Returns no output data.
  \item Reverts state changes made within the call (but the gas is still consumed).
\end{enumerate}

\subsection{Input Validation}

Every precompile validates inputs before processing:

\begin{lstlisting}[language=Go,caption={Input validation pattern}]
func (c *dilithiumVerify) Run(input []byte) ([]byte, error) {
    // Check total length
    if len(input) < PubKeyLen+SigLen+32 {
        return nil, errInvalidInputLength
    }

    // Decode public key
    pk, err := decodeDilithiumPubKey(input[:PubKeyLen])
    if err != nil {
        return nil, errInvalidPublicKey
    }

    // Decode signature
    sig, err := decodeDilithiumSig(input[PubKeyLen : PubKeyLen+SigLen])
    if err != nil {
        return nil, errInvalidSignature
    }

    // Message is the remainder
    msg := input[PubKeyLen+SigLen:]

    // Verify
    valid := dilithium.Verify(pk, msg, sig)

    output := make([]byte, 32)
    if valid {
        output[31] = 1
    }
    return output, nil
}
\end{lstlisting}

\subsection{Denial-of-Service Resistance}

The \texttt{RequiredGas} function must accurately bound the computation cost. For
variable-cost operations (e.g., pairing with $k$ pairs), the gas scales linearly:

\begin{lstlisting}[language=Go,caption={Variable gas cost for pairing}]
func (c *bls12381Pairing) RequiredGas(input []byte) uint64 {
    if len(input) == 0 || len(input)%128 != 0 {
        return 0 // Will fail in Run
    }
    k := uint64(len(input)) / 128
    return 43000 + k*70000 // Base + per-pair cost
}
\end{lstlisting}

%% ============================================================
